%==============================================================================
%  RÉSUMÉ / ABSTRACT
%==============================================================================
\chapter*{Résumé}
\addcontentsline{toc}{chapter}{Résumé}
\mtcaddchapter

Les plateformes d'apprentissage en ligne ouvertes et massives (MOOC) souffrent de
taux d'abandon élevés, largement imputables à une baisse d'engagement que les
dispositifs actuels peinent à détecter et à corriger. Ce mémoire présente
\emolearn{}, un \emph{prototype de recherche} qui articule trois briques : (i) la
\textbf{perception} de l'état d'engagement à partir de signaux multimodaux
(visage et voix), (ii) une \textbf{décision} pédagogique fondée sur une taxonomie
d'états d'engagement, et (iii) la \textbf{génération} automatique de contenu
pédagogique adaptatif au moyen d'un grand modèle de langue (LLM), complété par
un mécanisme de \textbf{récupération augmentée} (RAG) qui ancre la génération
sur une base documentaire pédagogique vectorisée. Le système
s'appuie sur le corpus d'engagement \textsc{DAiSEE} et sur le corpus émotionnel
\textsc{IEMOCAP}, et propose une chaîne complète de bout en bout dont chaque
maillon — y compris ses limites — est explicité et, lorsque c'est possible,
mesuré. Nous décrivons l'architecture, le modèle engagement\,$\rightarrow$\,adaptation,
les vingt patrons de génération, ainsi qu'un protocole d'évaluation multi-niveaux.

\vspace{0.5cm}
\noindent\textbf{Mots-clés :} apprentissage adaptatif, engagement de l'apprenant,
informatique affective, fusion multimodale, IA générative, RAG, MOOC, DAiSEE,
IEMOCAP.

\vspace{1cm}
\section*{Abstract}

Massive Open Online Courses (MOOCs) suffer from high dropout rates, largely driven
by a decline in learner engagement that current platforms struggle to detect and
remediate. This thesis introduces \emolearn{}, a research prototype that combines
three components: (i) \emph{perception} of the engagement state from multimodal
signals (face and voice), (ii) a pedagogical \emph{decision} grounded in an
engagement-state taxonomy, and (iii) automatic \emph{generation} of adaptive
learning content using a large language model (LLM), augmented by
\textbf{retrieval-augmented generation} (RAG) over a vectorized pedagogical
knowledge base. The system builds on the
\textsc{DAiSEE} engagement dataset and the \textsc{IEMOCAP} emotional corpus, and
proposes a full end-to-end pipeline whose every component --- including its
limitations --- is made explicit and, where feasible, measured.

\vspace{0.5cm}
\noindent\textbf{Keywords:} adaptive learning, learner engagement, affective
computing, multimodal fusion, generative AI, RAG, MOOC, DAiSEE, IEMOCAP.
